\item \textbf{[4 pts.] Thermal conductivity} \\
Using the mineralogical composition chart and the thermal conductivities
listed in the table below, compute the thermal conductivity for

\par\medskip
\begin{figure}[hbt]
  \centering{
    \includegraphics[angle=90]{figures/igneous_rock_minerals.eps}
  }
  \caption{Mineralogical compositions of igneous rocks.}
  \label{fig:geo}
\end{figure}
\medskip

\par\medskip
\begin{table}[hbt]
    \centering
    \begin{threeparttable}
    \caption{Thermal conductivity of common materials.}
    \label{tab:tc}
    \begin{tabular}{@{}lc@{}}
        \toprule
        {} &
        $k$ \\
        mineral &
        \si{W.m^{-1}.K^{-1}} \\
        \midrule
        $\alpha$-quartz           & 8.8 \\
        orthoclase                & 1.8 \\
        plagioclase, felsic       & 1.8 \\
        plagioclase, mafic        & 1.6 \\
        mica, muscovite           & 2.1 \\
        mica, biotite             & 1.9 \\
        hornblende (amphibole)    & 2.7 \\
        clinopyroxene             & 4.3 \\
        orthopyroxene             & 3.4 \\
        olivine                   & 4.8 \\
        spinel                    & 9 \\
        garnet                    & 3.9 \\
        \bottomrule
    \end{tabular}
    \begin{tablenotes}
        \item All conductivities estimated at \SI{25}{^\circ C}.
    \end{tablenotes}
    \end{threeparttable}
\end{table}
\medskip

\begin{enumerate}
    \item granite
    \item diorite
    \item gabbro, assuming all pyroxene is clinopyroxene.
    \item peridotite

    Phase changes associated with pressure turn plagioclase into spinel and
    then into garnet.  Therefore, compute the thermal conductivity of
    peridotite with a composition of 67\% olivine, 17\%
    orthopyroxene, 8\% clinopyroxene, and 8\% garnet.
\end{enumerate}
\FloatBarrier

